% ============================================================
% Formulario, argomento 24: minimi quadrati e stima ricorsiva (dispensa, capitolo 24)
% ============================================================
\argomento{Minimi quadrati e stima ricorsiva}{Dispensa: capitolo 24}

\voce{modellolineare}{Modello lineare}{$p$ vettori di $n$ caratteristiche e i $p$ valori misurati}
  {\theta(i) = \vett{a}^T \vett{x}(i) + \epsilon_i, \quad i = 1, \ldots, p \qquad \matr{X} = \left[ \vett{x}(1)\ \cdots\
   \vett{x}(p) \right] \in \R^{n \times p}}
  {\variante{le osservazioni}{\vett{y} = [\theta(1), \ldots, \theta(p)]^T \in \R^p}
   \variante{somma dei quadrati degli errori}{\norm{\vett{\epsilon}}^2 = \sum_{i=1}^{p} \left( \vett{a}^T \vett{x}(i)
     - \theta(i) \right)^2 = \norm{\matr{X}^T \vett{a} - \vett{y}}^2}
   \variante{norma}{\norm{\vett{v}}^2 = \vett{v}^T \vett{v} = \sum_i v_i^2}}

\voce{minimiquadrati}{Stima ai minimi quadrati (LS)}{\emph{Least Squares}; $\matr{X}\matr{X}^T$ invertibile}
  {\vett{a}_{\text{LS}} = \left( \matr{X} \matr{X}^T \right)^{-1} \matr{X}\, \vett{y}}
  {\variante{equazioni normali}{\matr{X}\matr{X}^T \vett{a} = \matr{X}\vett{y}}
   \variante{retta $\theta = a_1 + a_2 t$, con $\vett{x}(i) = [1, t_i]^T$}{\matr{X}\matr{X}^T = \begin{bmatrix} p &
     \sum t_i \\ \sum t_i & \sum t_i^2 \end{bmatrix}, \quad \matr{X}\vett{y} = \begin{bmatrix} \sum \theta_i \\ \sum
     t_i \theta_i \end{bmatrix}}
   \variante{servono $p \geq n$ dati e caratteristiche linearmente indipendenti}{}
   \variante{è lo stimatore lineare MMSE con le medie sui dati}{\frac{1}{p} \matr{X}\matr{X}^T \to \matr{R}, \quad
     \frac{1}{p} \matr{X}\vett{y} \to E[\Theta \vett{X}]}}

\voce{shermanmorrison}{Formula di Sherman-Morrison}{inversa dopo un aggiornamento di rango 1; $1 + \vett{v}^T
  \matr{R}^{-1} \vett{u} \neq 0$}
  {\left( \matr{R} + \vett{u}\vett{v}^T \right)^{-1} = \matr{R}^{-1} - \frac{\matr{R}^{-1} \vett{u}\, \vett{v}^T
   \matr{R}^{-1}}{1 + \vett{v}^T \matr{R}^{-1} \vett{u}}}
  {}

\voce{rls}{Minimi quadrati ricorsivi}{arriva il dato $\vett{x}(p + 1)$, $\theta(p + 1)$; in breve $\vett{x} =
  \vett{x}(p + 1)$}
  {\matr{R}(p + 1) = \matr{R}(p) + \vett{x}\vett{x}^T \qquad \matr{R}^{-1}(p + 1) = \matr{R}^{-1}(p) -
   \frac{\matr{R}^{-1}(p)\, \vett{x}\vett{x}^T \matr{R}^{-1}(p)}{1 + K}}
  {\variante{con il numero}{K = \vett{x}^T \matr{R}^{-1}(p)\, \vett{x}, \quad \matr{R}(p) = \matr{X}(p)\matr{X}^T(p)}
   \variante{la nuova stima}{\vett{a}(p + 1) = \left[ \matr{I}_n - \frac{\matr{R}^{-1}(p)\, \vett{x}\vett{x}^T}{1 + K}
     \right] \left[ \vett{a}(p) + \theta(p + 1)\, \matr{R}^{-1}(p)\, \vett{x} \right]}
   \variante{letta come correzione \oltre}{\vett{a}(p + 1) = \vett{a}(p) + \frac{\matr{R}^{-1}(p)\, \vett{x}}{1 + K}
     \left[ \theta(p + 1) - \vett{x}^T \vett{a}(p) \right]}
   \variante{costo $O(n^2)$ operazioni per ogni dato, indipendentemente da $p$}{}}
